\section{Related Work}
\label{submission:sec:related}

\textbf{Model Merging and Task Arithmetic.} Weight-space model merging has gained massive traction as an alternative to multi-task fine-tuning, as it allows combining specialized capabilities without retraining the shared backbone. Task Arithmetic \cite{ilharco2022editing} computes task-specific updates (task vectors) relative to a pre-trained base model, showing that linear operations in parameter space correspond to predictable semantic changes in behavior. However, conflicting signs or overlapping updates in dense task vectors often lead to destructive interference when merging multiple models. To address this, TIES-Merging \cite{yadav2023ties} filters out small-magnitude updates, consensus-votes on sign directions, and averages only sign-consistent updates. DARE \cite{yu2023language} randomly drops a fraction of delta parameters (using Bernoulli dropout) and rescales the remaining ones to keep expectations constant. While TIES and DARE incorporate weight pruning as a denoising step, they do so uniformly and do not explicitly optimize for post-hoc storage footprint or leverage optimizer-driven geometry.

\textbf{Loss Landscapes, Flatness, and SAM.} The geometry of the loss landscape is closely tied to a model's generalization capabilities \cite{foret2021sharpness}. Sharpness-Aware Minimization (SAM) explicitly optimizes for flat minima by perturbing parameters in the direction of the greatest gradient increase during training. Recently, \citet{saim2024deconstructing} systematically deconstructed sharpness-aware isotropic merging, discovering that optimizer-driven flatness (fine-tuning experts using SAM) is the primary causal driver of multi-task merging success, as it allows experts to reside in wide, overlapping valleys. In this work, we investigate this landscape geometric connection under post-hoc weight sparsification, revealing a counter-intuitive finding: while optimizer-driven flatness successfully stabilizes dense weight merging, it does not inherently buffer task vectors against unstructured, coordinate-aligned magnitude pruning. We provide an in-depth comparative analysis to dissect this phenomenon.

\textbf{Model Compression and Edge Sparsification.} Reducing the computational and storage footprint of deep networks is essential for wild deployments \cite{han2015learning, chmiel2020robust}. Traditional compression methods, such as magnitude pruning, quantization, and distillation, are typically applied to individual models post-training. However, compressing merged or multi-task networks remains relatively under-explored. In edge scenarios where bandwidth is constrained, sending fully dense parameters over-the-air is extremely expensive. By pruning task vectors globally to extreme levels (e.g., 90\% to 95\% sparsity), we enable storing delta updates in compressed sparse formats like Coordinate list (COO) or Compressed Sparse Row (CSR), reducing storage requirements by 10--20$\times$.

\textbf{Avoiding Test-Time Overfitting.} Recent works, such as AdaMerging, attempt to search for optimal layer-wise merging coefficients at test-time using small calibration sets. However, \citet{adamerging2024sanity} sanity-checked these adaptive techniques and uncovered the ``Overfitting-Optimizer Paradox'', showing that layer-wise coefficient tuning leads to severe transductive overfitting on the calibration set while simple flat merging generalizes much better. This finding strongly motivates our training-free, post-hoc offline pruning and merging strategy, which avoids fragile test-time optimization entirely and focuses on building robust parameter-space representations.
